\documentclass{article}

\usepackage{amsmath}
\usepackage{graphicx}

\title{Planetary Intelligence Quotient Network}
\author{Author}
\date{}

\begin{document}

\maketitle

\begin{abstract}
Technological civilization increasingly consists of integrated automated systems including artificial intelligence infrastructure, robotics, cloud computing, satellite sensing, and global logistics networks. These systems form a cybernetic network embedded within planetary infrastructure.

This paper introduces the Planetary Intelligence Quotient Network (PIQNet), a neural architecture designed to estimate the emergent intelligence level of technological civilization. The model combines graph neural networks to represent infrastructure integration and transformer-based temporal modeling to capture technological evolution. The resulting model outputs a scalar Planetary Intelligence Quotient representing the amplification of intelligence across the global technological network.
\end{abstract}

\section{Introduction}

Modern civilization increasingly behaves as a distributed cybernetic system composed of interacting technological subsystems. Understanding the intelligence level of this integrated infrastructure is an important research problem.

\section{Model}

The Planetary Intelligence Quotient is estimated using

\[
PIQ = f_\theta(G_t, T_t)
\]

where \(G_t\) represents infrastructure networks and \(T_t\) represents temporal technological evolution.

\section{Architecture}

The model combines

\begin{itemize}
\item Graph Neural Networks
\item Transformer Temporal Models
\item Regression Head for PIQ estimation
\end{itemize}

\section{Conclusion}

PIQNet provides a framework for analyzing the intelligence level of technological civilization using neural network models.

\end{document}
